\begin{table}[htbp]
\centering
\caption{Prespecified evaluation plan for the deterministic red-flag (SSNHL/neurologic) safety layer. Sensitivity (red-flag recall) is the primary metric because a missed red flag is the costly error; but a perfect score on a finite test set does not guarantee zero real-world misses, and this limitation is stated explicitly.}
\label{tab:redflag}
\small
\begin{tabular}{p{0.28\linewidth}p{0.66\linewidth}}
\toprule
Evaluation element & Prespecification \\
\midrule
Vignette set & $\geq$200 vignettes; composition fixed: real deidentified clinical cases, expert-authored cases, and synthetic cases in a stated ratio (e.g., 40/30/30\%). \\
Reference labeling & Each vignette independently labeled by $\geq$2 otolaryngology/audiology experts; disagreements adjudicated. \\
Inter-rater agreement & Reported (Cohen's/Fleiss' $\kappa$) before adjudication. \\
Primary metric & Sensitivity (red-flag recall) with 95\% CI. \\
Secondary metrics & Specificity and over-referral rate with 95\% CIs. \\
Subgroup performance & By laterality (unilateral/bilateral), onset (acute/gradual), and vertigo/neurologic signs. \\
Robustness & Negation, ambiguous temporal expressions, and colloquial patient phrasing explicitly tested. \\
Error analysis & Every missed red flag audited and reported as a critical failure. \\
Stated limitation & Finite-set sensitivity of 1.0 is a target, not a guarantee of zero misses in deployment. \\
\bottomrule
\end{tabular}
\end{table}
